% !TEX root = ../aa_SoM_Chapters.tex

% ö = \%c3\%b6
% ä = \%C3\%A4

\xdef\sectiontitle{\Isectiontitle} %aiheuttama

%%%%%%%%%%%%%%%
\begin{frame}[label=1_4_a]%[label=1_4_TOC1]
\frametitle{\texorpdfstring{\culine[\sectiontitleulinecolor]{\sectiontitle}}{\sectiontitle} }

\vspace{0.5cm}

\begin{itemize}[leftmargin=0.5cm,itemsep=2mm]
    \item[1.1]<1-> \hyperlink{1_1_a}{\IaSubsectiontitle}
    \item[1.2]<1-> \hyperlink{1_2_a}{\IbSubsectiontitle}	
    \item[1.3]<1-> \hyperlink{1_3_a}{\IcSubsectiontitle}	
    \item[1.4]<1-> \hyperlink{1_4_a}{\Alert[red]{\IdSubsectiontitle}}
    \item[1.5]<1-> \hyperlink{1_5_a}{\IeSubsectiontitle}
    \item[1.6]<1-> \hyperlink{1_6_a}{\IfSubsectiontitle}
    \item[1.7]<1-> \hyperlink{1_7_a}{\IgSubsectiontitle}
\end{itemize}

\vspace{-1cm}
\begin{itemize}[leftmargin=9.5cm,itemsep=2mm]
    \item[1.8]<1-> \hyperlink{1_8_a}{\IhSubsectiontitle}
	\item[1.9]<1-> \hyperlink{1_9_a}{\IiSubsectiontitle}
	\item[1.10]<1-> \hyperlink{1_10_a}{\IjSubsectiontitle}
	\item[1.11]<1-> \hyperlink{1_11_a}{\IkSubsectiontitle}
\end{itemize}

%\endofslide 
\end{frame}
%%%%%%%%%%%%%%%

%\xdef\IbSubsectiontitle{Entropy and Temperature}
\xdef\sectiontitle{\IdSubsectiontitle}

%%%%%%%%%%%%%%%
\begin{frame}%[label=1_4_a]
\frametitle{\texorpdfstring{\culine[\sectiontitleulinecolor]{\sectiontitle}}{\sectiontitle} }

\begin{itemize}
	\item Applying equation for $P_{n_{i}}$\appear{, shown \hyperlink{1_3_Pn}{on earlier slide}, is tedious.} \appear{For large systems, it is feasible to use \CDAlert[red]{Boltzmann probability}:}
\[ 
 P \textrm{((} E_{n} \textrm{))}  \equal \appear{A} \appear{e^{-E_{n} / k_{B} T} }
 \]
\appear{It is seen in Figure 6 that Boltzmann works fairly well} \appear{even for smaller systems like the one given}
\end{itemize}

\vspace{2.5mm}
% 6.5 cm = midway, 10 cm = 3/4 
\begin{minipage}{5.5cm}
\begin{itemize}
	\item For $P \textrm{((} E_{n} \textrm{))} $ to be a qualified probability function\appear{, it needs to be normalized}\appear{, i.\,e. we need to choose constant $A$ properly}
%\[
%\sum P \textrm{((} E_{n} \textrm{))} =\sum_{n} A e^{-E_{n} / k_{B} T}=1 \Rightarrow A=\frac{1}{\sum_{n} e^{-E_{n} / k_{B} T}} \Rightarrow P \textrm{((} E_{n} \textrm{))} =\frac{e^{-E_{n} / k_{B} T}}{\sum_{n} e^{-E_{n} / k_{B} T}}
%\[
%And thus, e.g. average energy can be obtained by
%\[
% < E >\, =\sum_{n} E_{n} P \textrm{((} E_{n} \textrm{))} =\sum_{n} E_{n} \frac{e^{-E_{n} / k_{B} T}}{\sum_{n} e^{-E_{n} / k_{B} T}}=\frac{\sum_{n} E_{n} e^{-E_{n} / k_{B} T}}{\sum_{n} e^{-E_{n} / k_{B} T}}
%\[
\end{itemize}
\end{minipage}

\xdef\loc{south east}
\begin{tikzpicture}[remember picture,overlay] % anchor=south
    \node (A) [yshift=-0.25*\figYshift,xshift=\figXshift, anchor=\loc] at (current page.\loc) {\includegraphics[page=1,width=8.2cm]{Figures_booklet/SOM_9_Statistical_mechanics_figures/SOM_Figure_6.png}}; %scale=\figscale. % 8cm max hei
\end{tikzpicture}

\endofslide 
%\endofsection
\end{frame}
%%%%%%%%%%%%%%%

%%%%%%%%%%%%%%%
\begin{frame}%[label=1_4_a]
\frametitle{\texorpdfstring{\culine[\sectiontitleulinecolor]{\sectiontitle}}{\sectiontitle} }

\begin{itemize}
\item[] \[
\begin{aligned}
\nsum \appear{P \textrm{((} E_{n} \textrm{))}} &  \equal  \appear{\nsum_{n}} \appear{A} \appear{e^{-E_{n} / k_{B} T} } \equal  \appear{1} \qquad \Rightarrow \qquad \appear{A}  \equal \frac{1}{\nsum_{n} e^{-E_{n} / k_{B} T}} \\
 \appear{\Rightarrow} \qquad \quad \appear{P \textrm{((} E_{n} \textrm{))}} & \equal \frac{e^{-E_{n} / k_{B} T}}{\nsum_{n} e^{-E_{n} / k_{B} T}}
\end{aligned}
\]
\item And thus, e.g. \CDAlert[red]{average energy $< E >\,$}$^*$ can be obtained by\footnote{{\appear{$^*$\;$< E >\,$ is the  \CDAlert[red]{expected value}{https://en.wikipedia.org/wiki/Expected_value}, defined as $\mathrm{E}(X) \Equiv \sum_{i=1}^n x_i p_i \,$.} \appear{Memorize~this.}}}
\[
\begin{aligned}
 < E >\, &  \equal \appear{\nsum_{n}} \appear{E_{n}} \appear{P \textrm{((} E_{n} \textrm{))}} \\
 & \equal \appear{\nsum_{n} E_{n}} \frac{e^{-E_{n} / k_{B} T}}{\nsum_{n} e^{-E_{n} / k_{B} T}} \\
 & \equal \fracc{\nsum_{n} E_{n} e^{-E_{n} / k_{B} T}}{\nsum_{n} e^{-E_{n} / k_{B} T}}
\end{aligned}
\]
\end{itemize}

%\xdef\loc{south east}
%\begin{tikzpicture}[remember picture,overlay] % anchor=south
%    \node (A) [yshift=-0.25*\figYshift,xshift=\figXshift, anchor=\loc] at (current page.\loc) {\includegraphics[page=1,width=8.2cm]{Figures_booklet/SOM_9_Statistical_mechanics_figures/SOM_Figure_6.png}}; %scale=\figscale. % 8cm max hei
%\end{tikzpicture}

\endofslide 
%\endofsection
\end{frame}
%%%%%%%%%%%%%%%

%%%%%%%%%%%%%%%
\begin{frame}
\frametitle{\texorpdfstring{\culine[\sectiontitleulinecolor]{\sectiontitle}}{\sectiontitle} }

\begin{itemize}
	\item For the Einstein's solid \appear{we had $E_{n}=n \hbar \omega_{0}$}\appear{, now let's plug this into above equation, and plot the Boltzmann distribution (Fig.~7)}
	\appear{\xdef\loc{south east}
\begin{tikzpicture}[remember picture,overlay] % anchor=south
    \node (A) [yshift=-0.25*\figYshift,xshift=\figXshift, anchor=\loc] at (current page.\loc) {\includegraphics[page=1,width=8cm]{Figures_booklet/SOM_9_Statistical_mechanics_figures/SOM_Figure_7.png}}; %scale=\figscale. % 8cm max hei
\end{tikzpicture}}

\item Variation of Boltzmann probability distribution $P \textrm{((} E_{n} \textrm{))}$ for 3 different temperature ranges \appear{(from $k_{\mathrm{B}} T\!<\!\hbar \omega_{0}$ to $k_{\mathrm{B}} T\!>\!\hbar \omega_{0}$):}
% 6.5 cm = midway, 10 cm = 3/4 
\end{itemize}

\begin{minipage}{6cm}
\begin{itemize}
	\item[] \begin{itemize}[itemsep=2mm]
	\item $T=\Frac{1}{5} \hbar \omega_{0} / k_{B}$
\item $T=\hbar \omega_{0} / k_{B}$
\item $T=5 \hbar \omega_{0} / k_{B}$
\end{itemize}

\item Obviously at lower temperatures it is very
	unlikely for the oscillators to possess
	several quanta \appear{(of $\hbar \omega_{0}$)}\appear{, i.\,e. to get
	to higher energy states}
\end{itemize}
\end{minipage}
\vspace{2.5mm}


\endofslide 
%\endofsection
\end{frame}
%%%%%%%%%%%%%%%

%%%%%%%%%%%%%%%
\begin{frame}
\frametitle{\texorpdfstring{\culine[\sectiontitleulinecolor]{\sectiontitle}}{\sectiontitle} }

%\begin{itemize}
%	% 6.5 cm = midway, 10 cm = 3/4 
%\item[]
\begin{minipage}{7cm}
\begin{itemize}
	\item We want to often know how
		many particles \appear{(oscillators)} \appear{are at
		certain state}
		\appear{(possess a certain number of quanta)}

\item Total number of particles is still $N$
\end{itemize}
\end{minipage}

\vspace{2.5mm}
\begin{itemize}

\item We define \CDAlert[red]{occupation number} $\mathbb{N}$: %$\mathcal{N}$
\[ 
 \mathbb{N} \textrm{((} E_{n} \textrm{))} _{B o l t z m a n n}  \equal  \appear{N P \textrm{((} E_{n} \textrm{))}}  \equal  \appear{N\,} \appear{A e^{-E_{n} / k_{B} T}} \appear{\,, \qquad  < E >\,  \equal \frac{\sum_{n} E_{n} \mathbb{N} \textrm{((} E_{n} \textrm{))} }{\sum_{n} \mathbb{N} \textrm{((} E_{n} \textrm{))} } }
 \]

\item Using occupation number $\mathbb{N}$\appear{, we get:}
\[ \bg[2.4]{\text{\bigsymbolsfont(}}\!\! 
\appear{{\footnotesize
\begin{array}{c}
\text {number of} \\ \text {particles of} \\ \text {energy E}
\end{array}  }}
 \!\!\bg[2.4]{\text{\bigsymbolsfont)}}   \equal  \appear{\bg[2.4]{\text{\bigsymbolsfont(}\! }
\appear{{\footnotesize
\begin{array}{c}
\mathbb{N} \text {, number of} \\ \text{particles in a given} \\ 
\text {state of energy E}
\end{array}}}
\!\!\bg[2.4]{\text{\bigsymbolsfont)}}  } \appear{\times} \appear{ \bg[2.4]{\text{\bigsymbolsfont(}}\!\!  
\appear{{\footnotesize 
\begin{array}{c}
\text {number of} \\ \text{states of} \\ \text {energy E}
\end{array}}}
\!\!\bg[2.4]{\text{\bigsymbolsfont)}} }
\]
\end{itemize}

\xdef\loc{north east}
\begin{tikzpicture}[remember picture,overlay] % anchor=south
    \node (A) [yshift=\figYshift,xshift=\figXshift, anchor=\loc] at (current page.\loc) {\includegraphics[page=1,width=6.5cm]{Figures_booklet/SOM_9_Statistical_mechanics_figures/SOM_Figure_7.png}}; %scale=\figscale. % 8cm max hei
\end{tikzpicture}

\endofslide 
%\endofsection
\end{frame}
%%%%%%%%%%%%%%%

%%%%%%%%%%%%%%%
\begin{frame}
\frametitle{\texorpdfstring{\culine[\sectiontitleulinecolor]{\sectiontitle}}{\sectiontitle} }

\begin{itemize}
	\item Now we can write: 
	\[
	< E >\,  \equal \fracc{\sum_{n} E_{n} \mathbb{N} \textrm{((} E_{n} \textrm{))} }{\sum_{n} \mathbb{N} \textrm{((} E_{n} \textrm{))} }
	\]
\item Often the quantum energy levels are so close to each other that the \CDAlert[red]{sums can be replaced by integrals}

\item Generally, the ratio between probabilities of two neighbouring energy levels is
\[ 
 \Frac{P \textrm{((} E_{n} \textrm{))} }{P \textrm{((} E_{n+1} \textrm{))} }  \equal \fracc{e^{-E_{n} / k_{B} T}}{e^{-E_{n+1} / k_{B} T}} \equal \appear{e^{ \textrm{((} E_{n+1}-E_{n} \textrm{))}  / k_{B} T} }
 \]
\item So if $E_{n+1}-E_{n}$ is much smaller than $k_{B} T$\appear{, then the ratio really is close to 1}\appear{, and the sum can be transformed into an integral}

\end{itemize}

\endofslide 
%\endofsection
\end{frame}
%%%%%%%%%%%%%%%


%%%%%%%%%%%%%%%
\begin{frame}
\frametitle{\texorpdfstring{\culine[\sectiontitleulinecolor]{\sectiontitle}}{\sectiontitle} }

\begin{itemize}
	\item This could be seemingly performed by taking the limit
\[ 
  < E >\,  \equal \frac{\sum_{n} E_{n} \mathbb{N} \textrm{((} E_{n} \textrm{))} }{\sum_{n} \mathbb{N} \textrm{((} E_{n} \textrm{))} } \qquad \appear{\Rightarrow} \qquad \frac{\int_{n} E \mathbb{N}(E) d n}{\int_{n} \mathbb{N}(E) d n} 
 \]
	\item However, we need to be careful when transforming sums into continuous integrals
	
\item First of all, the quantum number $n$ is not necessarily the most convenient integration variable here\appear{, energy is a better quantity}

\item Therefore we should have:
\[ 
  < E >\,  \equal \frac{\int E \mathbb{N}(E) d n}{\int \mathbb{N}(E) d n} \equal \frac{\int E \mathbb{N}(E) \frac{d n}{d E} \appear{d E}}{\int \mathbb{N}(E) \appear{\Frac{d n}{d E} d E}} \equal \frac{\int E \mathbb{N}(E) D(E) d E}{\int \mathbb{N}(E) D(E) d E} 
 \]

\end{itemize}

\endofslide 
%\endofsection
\end{frame}
%%%%%%%%%%%%%%%

%%%%%%%%%%%%%%%
\begin{frame}
\frametitle{\texorpdfstring{\culine[\sectiontitleulinecolor]{\sectiontitle}}{\sectiontitle} }

\begin{itemize}
\item[] %Therefore we should have:
\[ 
  < E >\,  \equal \Frac{\int E \mathbb{N}(E) d n}{\int \mathbb{N}(E) d n}  \equal \Frac{\int E \mathbb{N}(E) \Frac{d n}{d E} d E}{\int \mathbb{N}(E) \Frac{d n}{d E} d E} \equal \Frac{\int E \mathbb{N}(E) D(E) d E}{\int \mathbb{N}(E) D(E) d E} 
 \]
\item Here we have introduced a quantity known as the \CDAlert[fuchsia]{density of} \CDAlert[fuchsia]{states $D(E)$}\appear{, indicating how closely the energy levels are packed:}
\[ 
 D(E) \equiv \fracc{\text{\small differential number of states within range dE close to E}}{d E} \equal \fracc{d n}{d E} 
 \]

%\fontsize{13}{15}\selectfont 
\item In practise, most (if not all) calculations of mean values \appear{are calculated as integrals}
\implies[\normal]{ \CDAlert[red]{Density of states $D$ (or DOS)} \appear{\CDAlert[red]{is a central and omnipresent}} \\ \appear{\CDAlert[red]{concept in solid-state physics!}} }
\implies[\normal]{ \Alert[tunired]{Re-introduced in our course on solid-state physics..} }
\end{itemize}

\endofslide 
%\endofsection
\end{frame}
%%%%%%%%%%%%%%%


%%%%%%%%%%%%%%%
\begin{frame}
\frametitle{\texorpdfstring{\culine[\sectiontitleulinecolor]{\sectiontitle}}{\sectiontitle} }

\begin{itemize}
\item Harmonic oscillator provides an easy example to demonstrate \appear{how the density of states $D(E)$ can be calculated:} 
\[
\begin{aligned}
E_{n}  \equal \appear{n \hbar \omega_{0}} \qquad &\Rightarrow& \qquad
\frac{d n}{d E} &\equal \frac{\text{differential number of states}}{d E} \\
&\Rightarrow& \qquad \appear{d n} &  \equal \frac{1}{\hbar \omega_{0}} \appear{d E} \\
&\Rightarrow &    \frac{d n}{d E} & \equal \frac{1}{\hbar \omega_{0}}  \equal \appear{D(E)}
 \end{aligned}
\]

\item Admittedly, this example is somewhat boring\appear{/trivial}\appear{, as it gives a constant density of states}\appear{, but at least we can verify that the result makes sense} \appear{(energy states of harmonic oscillator are equispaced)}
\end{itemize}

%\xdef\loc{south west}
%\begin{tikzpicture}[remember picture,overlay] % anchor=south
%    \node (A) [yshift=2cm,xshift=-\figXshift, anchor=\loc] at (current page.\loc) {\includegraphics[page=1,width=6.8cm]{Figures_booklet/SOM_9_Statistical_mechanics_figures/SOM_Figure_4_cc.png}}; %scale=\figscale. % 8cm max hei
%\end{tikzpicture}

\endofslide 
%\endofsection
\end{frame}
%%%%%%%%%%%%%%%

%%%%%%%%%%%%%%%
\begin{frame}
\frametitle{\texorpdfstring{\culine[\sectiontitleulinecolor]{\sectiontitle}}{\sectiontitle} }

\begin{itemize}
	\item For a set of identical (but distinguishable) harmonic oscillators\appear{, we obtain the average energy:}
\[ 
  < E >\,  \equal \fracc{\nint E \mathbb{N}(E) D(E) d E}{\nint \mathbb{N}(E) D(E) d E}  \equal \frac{\nint E N A e^{-E_{n} / k_{B} T} \frac{1}{\hbar \omega_{0}} d E}{\nint N A e^{-E_{n} / k_{B} T} \frac{1}{\hbar \omega_{0}} d E}  \equal \appear{\ldots} \equal \appear{k_{B} T} 
 \]
\item Previously we obtained a result:
\[ 
  < E >\,  \equal \frac{\hbar \omega_{0}}{e^{\hbar \omega_{0} / k_{B} T}-1} 
 \]
\item If we have $\hbar \omega_{0} \ll k_{B} T$\appear{, then $\appear{x=}\frac{\hbar \omega_{0}}{k_{B} T} \appear{\ll 1}$}\appear{, and using} $\appear{e^{x}} \approx \appear{1+x}$\appear{, we obtain $ \appear{< E >} \approx \appear{k_{B} T}$}
\end{itemize}

%\endofslide 
\endofsection
\end{frame}
%%%%%%%%%%%%%%%

